\section{Related Work}
\label{sec:related}

\textbf{Unsupervised video object segmentation (VOS)}~\cite{ytvos2018, davis16, segtrackv2, FBMS59, moca} targets identifying and segmenting moving objects as foreground elements in video sequences, generating pixel-level binary masks that distinguish these objects from background content without any human-annotated labels. The task applies to both scenarios, including those with a single object instance and those with multiple concurrent instances.
It is worth noting that some materials~\cite{lee2024guided, matnet, ren2021reciprocal, yang2021learning} refer to video salient object detection (SOD) as \textit{Unsupervised VOS}, which is a different task that requires human-annotated labels. Current unsupervised video segmentation methods predominantly rely on detection based on motion or consistency cues. Motion Group~\cite{motiongroup} uses optical flow to train an unsupervised video object segmentation network, and it uses only optical flow as input. Similar to Motion Group~\cite{motiongroup}, many methods~\cite{emdriven, bootstrapping, guessmove, yang2019unsupervised, oclr, ventura2019rvos} use optical flow as input or use it to supervise the model. However, they all rely on optical flow estimators (e.g., RAFT~\cite{raft}) that are trained on human-annotated datasets, and most of them~\cite{emdriven, bootstrapping, guessmove, yang2019unsupervised} can only track one object at a time. While some unsupervised video object segmentation methods~\cite{oclr, ventura2019rvos} support multi-object tracking, their multi-object tracking performance is very limited, which makes them unsuitable for the video instance segmentation task.

\noindent\textbf{Unsupervised video instance segmentation (VIS)}~\cite{ytvis2019, wu2022seqformer, lin2020video, heo2022vita, wang2021end, zhang2023dvis} targets identifying and segmenting moving objects from backgrounds and discriminating between distinct instances, all without any human-annotated labels.
VideoCutLER~\cite{videocutler} is the first unsupervised video segmentation method that is built for the VIS task, without using any human-annotated labels in the whole pipeline. VideoCutLER~\cite{videocutler} establishes a milestone in unsupervised video instance segmentation (VIS), surpassing previous benchmarks by demonstrating multi-instance segmentation without using optical flow. While its key innovation --- a synthetic video generation pipeline leveraging spatial augmentations of CutLER~\cite{cutler} pseudo-labels (originally derived from ImageNet~\cite{imagenet}) --- enables remarkable performance, the approach remains fundamentally limited by synthetic-to-real domain discrepancies and rigid instance modeling. Specifically, the static object configurations in synthesized videos fail to capture authentic motion dynamics observed in natural video sequences.

\noindent\textbf{Looped self-training method} is the method that uses the pseudo labels generated by the model to retrain the model~\cite{crst, zou2018unsupervised, noisy_student, st3d, twist, cutler}. This method is widely used in different areas related to unsupervised or semi-supervised training, for example, unsupervised domain adaptation~\cite{crst, zou2018unsupervised, st3d}, semi-supervised image classification~\cite{noisy_student}, semi-supervised 3D instance segmentation~\cite{twist}, unsupervised object detection and segmentation~\cite{cutler}, and human-in-the-loop systems for object detection~\cite{kuznetsova2020open,papadopoulos2016we}, instance segmentation~\cite{papadopoulos2021scaling,delatolas2024learning}, and
3D shape segmentation~\cite{yi2016scalable}.
